\documentclass[../HDF5_RFC.tex]{subfiles}

\begin{document}

%% ======================================================================
%%  Section 12 – Resolved Design Decisions
%% ======================================================================
\section{Resolved Design Decisions}
\label{sec:resolved}

The following items were raised during review and have been resolved.

\verysubsection{R1}
API naming and signatures. The introspection API went through several
iterations: (1) a single \texttt{ssize\_t}-returning function with two output buffers
and ambiguous return semantics; (2) two \texttt{ssize\_t}-returning functions; (3) a
single merged \texttt{herr\_t} function with nine parameters. The final design uses
two focused \texttt{herr\_t}-returning functions---
\texttt{H5Pget\_filter\_name\_by\_idx} (6~params) and \texttt{H5Pget\_filter\_params\_by\_idx}
(5~params)---each with an explicit \texttt{size\_t*} output for string length. This
follows the newer HDF5 API convention of reporting sizes via output
parameters, keeps each signature clean, and avoids a nine-parameter monolith.
The \texttt{filter\_idx} parameter uses \texttt{unsigned} to match \texttt{H5Pget\_filter2}. See
\SectionRef{sec:introspection}.

\verysubsection{R2}
Memory bounds for \texttt{cd\_values}. The \texttt{set\_config} callback supports a
two-pass size query pattern: the library first calls with \texttt{cd\_values=NULL} to
get the required count, allocates accordingly, and calls again. See
\SectionRef{sec:callback-specs}.

\verysubsection{R3}
Raw fallback syntax in the string API. The library does NOT implement a
special \texttt{cd\_values=...} pass-through syntax within
\texttt{H5Pset\_filter\_by\_name}. Users who need to pass raw \texttt{cd\_values} to a v2
plugin should use the existing \texttt{H5Pset\_filter} API. This keeps the string
API clean.

\verysubsection{R4}
Pipeline-in-a-string as primary API. The single-string pipeline
variant (\SectionRef{sec:pipeline-string}) is NOT the primary API. The primary
API is \texttt{H5Pset\_filter\_by\_name} (one filter per call, appending), consistent
with established HDF5 conventions.

\verysubsection{R5}
DCPL string storage. The DCPL does NOT retain the original parameter
string. After resolution, only the standard \texttt{cd\_values}, filter ID, and flags
are stored. Introspection functions rely on \texttt{get\_config} to reconstruct a
human-readable string from \texttt{cd\_values}.

\verysubsection{R6}
VOL connector interaction. The string API resolves to \texttt{cd\_values}
before the VOL connector sees the DCPL. No changes to the VOL interface are
required. See \SectionRef{sec:thread-safety}.

\verysubsection{R7}
Name-based plugin discovery. The plugin loading infrastructure (H5PL)
is extended by giving filters the same discriminated-union key structure that
VFDs and VOLs already use (\texttt{H5PL\_filter\_key\_t} with a
\texttt{BY\_NAME}/\texttt{BY\_ID} kind field), and adding an
\texttt{H5Z\_check\_plugin\_load()} callback. This follows the existing VFD/VOL
pattern exactly---no new plugin type constants, no naming conventions, no
manifest files. The name registry (\SectionRef{sec:name-registry}) serves as a
cache after first discovery. Only v3 plugins are eligible for name-based
matching; v2 plugins are skipped. See \SectionRef{sec:name-discovery}.

\verysubsection{R8}
Version numbering for \texttt{H5Z\_class3\_t}. The new struct carries
\texttt{version = 2}. The dispatch logic in \texttt{H5Zregister()} is extended to a
three-way branch (two-way under \texttt{H5\_NO\_DEPRECATED\_SYMBOLS}). The
\texttt{H5Z\_class\_t\_vers} compatibility macro in \texttt{H5version.h} retains its default
of~2 (mapping to \texttt{H5Z\_class2\_t}); it is NOT bumped to~3 to avoid silent
struct-size mismatches. See \SectionRef{sec:version-numbering}.

\verysubsection{R9}
Built-in filter validation. \texttt{set\_config} callbacks for built-in
filters with non-trivial validation (SZIP, Scale-Offset) delegate to existing
internal validation code rather than reimplementing it. See
\SectionRef{sec:builtin-filters}.

\verysubsection{R10}
Error code for unsupported string config. The error
\texttt{H5E\_UNSUPPORTED} (not a new code) is used when a filter does not implement
\texttt{set\_config}, rather than a novel \texttt{H5E\_CANTENCODE} code. See
\SectionRef{sec:set-filter-by-name}.

\verysubsection{R11}
\texttt{H5Zlookup\_by\_name} does not trigger plugin loading. It is a pure
in-memory registry lookup with no I/O side effects. Only
\texttt{H5Pset\_filter\_by\_name} and \texttt{H5Zfilter\_avail\_by\_name} trigger on-demand
plugin search. See \SectionRef{sec:explicit-registration}.

\verysubsection{R12}
Introspection fallback for filters without \texttt{get\_config}.
\texttt{H5Pget\_filter\_params\_by\_idx} formats the raw \texttt{cd\_values} as a fallback
string using colon separators (e.g., \texttt{"cd\_values=1:0:0:1074921472"}) so that
introspection always yields something useful. The colon separator avoids
ambiguity with the comma-separated grammar and keeps the fallback string valid
under the grammar as a single key=value pair. See \SectionRef{sec:introspection}.

\verysubsection{R13}
Multiple names per filter ID (aliases). A single filter ID MAY have
multiple registered names. \texttt{H5Zregister\_name} allows alias registration.
Introspection returns the first-registered name for that ID. See
\SectionRef{sec:explicit-registration}.

\verysubsection{R14}
Numeric name parsing order. When the \texttt{name} argument to
\texttt{H5Pset\_filter\_by\_name} is a numeric string, the name registry is checked
first. Only if the name is not found is it parsed as a numeric filter ID. A
registered name always takes precedence. See \SectionRef{sec:set-filter-by-name}.

\verysubsection{R15}
\texttt{H5Z\_CLASS\_T\_VERS} and \texttt{H5Z\_class\_t\_vers} must not change. The
existing constant \texttt{H5Z\_CLASS\_T\_VERS} (= 1) is left unchanged. A new constant
\texttt{H5Z\_CLASS3\_T\_VERS} (= 2) is introduced for v3 structs. The
\texttt{H5Z\_class\_t\_vers} compatibility macro in \texttt{H5version.h} retains its default
value of~2 (mapping \texttt{H5Z\_class\_t} $\rightarrow$ \texttt{H5Z\_class2\_t}) and is NOT
bumped to~3. See \SectionRef{sec:version-numbering}.

\verysubsection{R16}
Developer helper for plugin authors. The library provides
\texttt{H5Zconfig\_get\_param} in \texttt{H5Zdevelop.h}, returning \texttt{htri\_t} with an
explicit \texttt{size\_t *val\_len} output. See \SectionRef{sec:config-get-param}.

\verysubsection{R17}
Filter availability by name. A new function
\texttt{H5Zfilter\_avail\_by\_name()} is provided as the by-name equivalent of
\texttt{H5Zfilter\_avail()}. Unlike \texttt{H5Zlookup\_by\_name} (pure registry lookup), it
triggers on-demand plugin loading. See \SectionRef{sec:explicit-registration}.

\verysubsection{R18}
Duplicate filter replacement semantics. \texttt{H5Pset\_filter\_by\_name}
follows the same replacement behavior as \texttt{H5Pset\_filter}: if a filter with
the same ID is already in the pipeline, it is replaced, not duplicated. The
two APIs are fully interchangeable on the same DCPL. See
\SectionRef{sec:set-filter-by-name}.

\verysubsection{R19}
Grammar forward compatibility. Future extensions to the parameter
string grammar MUST be backward-compatible: any string valid under the current
grammar MUST remain valid with the same semantics under future grammars. See
\SectionRef{sec:risks}.

\verysubsection{R20}
Parameterless filter rejection. \texttt{set\_config} callbacks for filters
that take no parameters (shuffle, fletcher32, nbit) SHOULD reject
non-\texttt{NULL}, non-empty \texttt{params} strings with \texttt{H5E\_BADVALUE}. See
\SectionRef{sec:set-filter-by-name}.

\verysubsection{R21}
\texttt{get\_config} round-tripping. The output of \texttt{get\_config} MUST conform
to the same grammar as \texttt{set\_config} input (including bare keys and quoted
values where appropriate) so that \texttt{get\_config} output is valid
\texttt{set\_config} input. See \SectionRef{sec:callback-specs}.

\verysubsection{R22}
Introspection return convention. The introspection functions
\texttt{H5Pget\_filter\_name\_by\_idx} and \texttt{H5Pget\_filter\_params\_by\_idx} return
\texttt{herr\_t} with explicit \texttt{size\_t*} output parameters for reporting string
lengths, following the newer HDF5 API convention. See
\SectionRef{sec:introspection}.

\verysubsection{R24}
Two-pass \texttt{cd\_nelmts} identity contract. The \texttt{set\_config} callback MUST
return the same \texttt{cd\_nelmts} on both passes of the size-query protocol. If
the second-pass value differs from the first (in either direction), the
library MUST return an error and MUST NOT write to the \texttt{cd\_values} buffer
with the inconsistent count. A smaller second-pass value would silently
truncate the populated array; a larger value risks a buffer overflow. Both
cases indicate a buggy callback. The RFC previously specified only the
overflow guard (second-pass exceeds allocated capacity). This resolved item
adds the symmetric lower-bound check and makes identical counts the explicit
contract. See \SectionRef{sec:callback-specs}.

\verysubsection{R25}
Empty keys and explicit-empty-value distinction. The grammar already
forbids empty keys (\textit{key} requires one or more characters), but this
constraint is now stated explicitly in the Rules section to guide
implementors. In addition, the distinction between a bare key (\texttt{"flag"})
and a key with an explicit empty value (\texttt{"key="}) is specified: both
produce an empty string delivered to the callback, but via different
syntactic paths. Callbacks MAY treat them identically in practice, but the
distinction is preserved at the protocol level. The string \texttt{"=value"} MUST
be rejected with \texttt{H5E\_BADVALUE}. See \SectionRef{sec:param-rules}.

\verysubsection{R26}
\texttt{h5repack} \texttt{UD=} disambiguation. \texttt{UD} is a forbidden canonical
filter name and MUST be rejected by the name registry. The \texttt{h5repack}
parser MUST always treat \texttt{UD=} as the legacy \texttt{UD=id,flags,n,v1,...}
syntax; disambiguation is handled by \texttt{h5repack} itself before calling
\texttt{H5Pset\_filter\_by\_name}, not by the library parser. Filters that
previously used \texttt{UD} as a canonical name MUST adopt a different name.
See \SectionRef{sec:cli-tools}.

\verysubsection{R27}
\texttt{h5dump} \texttt{PARAMS\_STRING} format change. The addition of
\texttt{PARAMS\_STRING} to \texttt{h5dump} output constitutes a format change that may
break downstream parsers. The field is gated behind a command-line option
(\texttt{--filter-params}) so that the default output is unchanged. The output
format when the option is active is specified normatively in the CLI Tools
section. See \SectionRef{sec:cli-tools}.

\verysubsection{R23}
Flags-override detection. A \texttt{set\_config} callback MAY modify the caller's
\texttt{flags} value; the modified value is written into the DCPL. There is no
warning mechanism for this: \texttt{H5E\_WARN} does not exist in HDF5, and pushing
to the error stack is ineffective because the stack is cleared on success.
Callers that need to detect a flag override can compare the flags they passed
with the flags returned by \texttt{H5Pget\_filter2} after
\texttt{H5Pset\_filter\_by\_name} returns. This preserves the existing SZIP behavior
of unconditionally forcing \texttt{H5Z\_FLAG\_OPTIONAL} without introducing a
spurious error-stack contract. See \SectionRef{sec:callback-specs} and
\SectionRef{sec:builtin-filters}.

\verysubsection{R28}
Normative string and parameter limits. The maximum parameter string length
(4096~bytes, \texttt{H5Z\_CONFIG\_STRING\_MAX}) and the maximum number of
parameters per string (64, \texttt{H5Z\_CONFIG\_MAX\_PARAMS}) are normative
MUST-level limits enforced by the library, not recommendations. Both
constants are public. See \SectionRef{sec:param-rules}.

\verysubsection{R29}
Locale independence for both input and output. The library's parameter
string parser, \texttt{set\_config} callbacks, and \texttt{get\_config} callbacks
MUST all operate in C-locale semantics. Process-wide locale settings MUST NOT
affect parsing or formatting. This ensures round-trip correctness in
environments with non-ASCII-compatible locale settings. See
\SectionRef{sec:param-rules} and \SectionRef{sec:callback-specs}.

\verysubsection{R30}
Canonical \texttt{set\_config} implementation template. A normative code
template is provided in the architecture section to guide plugin authors in
correctly implementing the two-pass \texttt{set\_config} protocol. All built-in
filter \texttt{set\_config} callbacks follow this structure. See
\SectionRef{sec:set-config-template}.

\verysubsection{R31}
\texttt{description} field in \texttt{H5Z\_class3\_t}. A new optional
\texttt{description} field preserves the display role of the v2 \texttt{name}
field. For v3 plugins, \texttt{H5Pget\_filter2} returns \texttt{description}
when non-\texttt{NULL}, falling back to \texttt{name} otherwise. This mitigates
the semantic change in the \texttt{name} field by allowing plugin authors to
provide backward-compatible human-readable text. See
\SectionRef{sec:class3-struct} and \SectionRef{sec:name-semantic-change}.

\verysubsection{R32}
Strict-aliasing safety in \texttt{H5Zregister()}. The function signature of
\texttt{H5Zregister()} accepts \texttt{const void~*cls}. To read the \texttt{version} field and
dispatch to the correct struct type, the implementation MUST NOT cast
\texttt{cls} directly to \texttt{const H5Z\_class2\_t~*} (the default \texttt{H5Z\_class\_t}).
If the caller passed an \texttt{H5Z\_class3\_t~*}, such a cast aliases one struct
type through a pointer to a different struct type, which is undefined behavior
under the C strict-aliasing rules even though both structs share \texttt{int
version} as their first member (the common-initial-sequence guarantee applies
only to union members, not to raw pointer casts). Instead, the version field
MUST be extracted via \texttt{memcpy} into a local variable before any struct cast
is performed:

\begin{minted}{c}
int version;
memcpy(&version, cls, sizeof(int));
if      (version == 1)   { const H5Z_class2_t *c2 = cls; ... }
else if (version == 2)   { const H5Z_class3_t *c3 = cls; ... }
\end{minted}

Casting through \texttt{char~*} or \texttt{unsigned char~*} is equally acceptable, as
the C standard explicitly permits aliasing through character types. See
\SectionRef{sec:version-numbering}.

\end{document}
